\documentclass{article}
\begin{document}


\section{Observational Analysis}

The raw data were downloaded from the official ExoALMA database without any additional calibration or reduction.
(Include details such as the observational time, right ascension and declination, how the ExoALMA team combined the different observations, and the data span.)
The dataset includes both the full field of view (10 arcseconds — to be double-checked) and the trimmed data (5 arcseconds), spanning Briggs weighting indices from 2.0 to 2.0.
This allows a thorough investigation of the trade-off between sensitivity and resolution in the final detection.

For the inner disk regions, the analysis was performed on residual images obtained by subtracting the best-fit disk model (a Gaussian model from frank) from the observations.
For the outer disk regions, the analysis was conducted on the cleaned images without any model subtraction, as the disk emission there is negligible.
Only the full-FOV image with a Briggs index of 2.0 was available, providing the highest sensitivity.

GoFish was used to construct the deprojected radial profile of the disk residual emission, taking into account the disk inclination and position angle.
All adopted protoplanetary disk parameters are listed in Table X.
The default bin sampling size of one pixel was used, assuming no beam correlation.
The median absolute deviation was used to estimate the uncertainty in each annulus, as it is less sensitive to outliers than the standard deviation.

The noise profile was then reprojected into the image plane to construct a two-dimensional signal-to-noise (S/N) map, from which high-S/N sources were identified.
A minimum 5 $\sigma$ threshold was required for a robust detection.
Although a circumplanetary disk (CPD) typically has a small spatial extent of only a few pixels, its light is spread over the entire beam because it is unresolved; therefore, the full beam area is required to reach such a high S/N.
With the sources catalogued, further visual inspection was performed to rule out false positives, such as bright points that are part of symmetric artefacts or background sources.







\section{Modelling Approach}

The modelling of the CPD emission was performed following Zhu et al (2018) and Andreews et al (2021). The formulae would be included here .

\section{Results and Discussion}
 Explain the optically thin thick regimes  clearly with formulae approximation. 


\end{document}